\documentclass[11pt]{article}
\usepackage[margin=1in]{geometry}
\usepackage{amsmath}
\usepackage{booktabs}
\usepackage[hidelinks]{hyperref}

\newcommand{\PEGSPending}[1]{\textbf{[PENDING: #1]}}

\title{Detectability and Warning Value of Prompt Elasto-Gravity Signals from Manila Trench Megathrust Earthquakes}
\author{\PEGSPending{author list, station operators, and warning partners}}
\date{Methods-first draft; not for submission}

\begin{document}
\maketitle

\begin{abstract}
We evaluate when prompt elasto-gravity signals (PEGS) from Manila Trench
megathrust scenarios can be detected and used for reliable earthquake
characterization under realistic South China and surrounding station geometry,
non-Gaussian continuous noise, high-microseism and typhoon conditions, station
outages, and explicit operational false-alarm constraints. The study compares
single-station and coherent physical baselines with a conditional graph model,
and quantifies decision-time gain relative to P-wave, W-phase, and GNSS
characterization using scenario- and location-specific tsunami arrivals.
\PEGSPending{Insert registered simulator, station, noise, baseline, and holdout
results.}
\end{abstract}

\section{Introduction}
PEGS theory, empirical-noise graph models, and a simulated real-time tsunami-
warning implementation already exist. The question addressed here is narrower:
whether South China and surrounding networks can provide reliable Manila Trench
decisions at a fixed operational false-alarm rate, and what station/noise
requirements define the feasible boundary.

\section{Scenario and station data}
Each source record binds segment, Mw, strike/dip/rake, depth, rupture dimensions,
slip, rise time, rupture velocity, and publication. Every tsunami arrival is
nested under that source and has a location and arrival definition. Qualitative
statements such as ``more than about three hours'' are not converted into exact
universal arrival times.

Station candidates require a contemporaneous BH/LH three-component epoch, full
response, waveform availability, licence, latency class, and noise quality.
Channel metadata or scalar sensitivity alone is insufficient. Noise windows are
frozen by season, day/night, calm sea, typhoon/microseism, urban/remote setting,
and outage state before model fitting.

\section{Methods}
\subsection{PEGS simulation benchmark}
The selected simulator must reproduce a published event at specified stations
and components, with predeclared waveform correlation, amplitude, and onset
tolerances, before Manila scenarios are generated. Earth model, reference frame,
source-time function, units, sampling, and code commit are recorded.

\subsection{Splits and domain shift}
Training, validation, and test sets separate rupture segments and noise dates.
Domain-out sources and stations are retained. Simultaneous station noise windows
preserve spatial correlation. Synthetic white noise is an ablation, not the
primary evaluation.

\subsection{Physical baselines}
Baselines proceed from single-station energy to explicit phase/time-aligned
multi-station coherent stacking, template/likelihood detection, and rapid source
parameter inversion. Station weights are explicit and L2-normalized. Fixed 0\%,
20\%, and 40\% outage masks retain exact station identities.

\subsection{False alarms and detection probability}
Noise-only scores calibrate a threshold in false alarms per 30-day month. If the
record is too short to resolve one exceedance at the target rate, the threshold
is placed above the observed maximum and the resolution limit is reported.
Held-out detection probability is then evaluated versus time. Earliest reliable
detection may require multiple consecutive qualifying time points.

\subsection{Reliable magnitude and warning value}
At every time we report magnitude MAE and bias, sensitivity for the predeclared
high-risk class (nominally $M_w\geq8.2$), lower-risk specificity, and prediction-
interval coverage. Reliable magnitude time is the first sustained point meeting
all criteria.

For location $x$,
\begin{align}
 \Delta t_{\mathrm{characterization}} &=
 t_{\mathrm{conventional,reliable\ Mw}}-t_{\mathrm{PEGS,reliable\ Mw}},\\
 T_{\mathrm{lead}}(x) &= t_{\mathrm{tsunami\ arrival}}(x)
 -t_{\mathrm{PEGS,reliable\ Mw}}.
\end{align}
Negative characterization gain is retained as a negative result.

\subsection{Network optimization and conditional GNN}
Existing, idealized, and incrementally augmented networks are compared by
detection time/probability, false alarms, magnitude error, robustness, and
station count. A graph model is trained only if held-out real-noise physical
baselines demonstrate stable information. It must outperform those baselines
under identical splits and thresholds, not merely fit synthetic waveforms.

\section{Results}
\PEGSPending{No validated PEGS simulator, frozen station inventory, or empirical
noise library is present. Do not report the unit-test fixtures as detection or
warning results.}

\section{Discussion}
A high-value negative result is an explicit station-count, placement, noise, and
magnitude boundary. Preceding tsunami arrival alone does not establish decision
value; the system must characterize hazard reliably and rarely enough to be
operationally usable.

\section{Data and code availability}
Station metadata, waveform licences, exact noise dates, simulator/code versions,
scenario provenance, splits, thresholds, compact metrics and checksums will be
recorded. Restricted waveforms and large synthetic corpora remain outside Git.

\section{Limitations}
\PEGSPending{Update after simulator benchmark, station/noise gate, physical
baselines, network optimization, and conditional graph ablations.}

\bibliographystyle{plain}
\bibliography{references}
\end{document}
